\chapter*{Kaynakça}
\addcontentsline{toc}{chapter}{Kaynakça}
\markboth{Kaynakça}{Kaynakça}
\fancyhead[RO]{\sffamily\scriptsize\color{gray}KAYNAKÇA}

Bu kaynakça, kitap boyunca kullanılan yazar--yıl atıflarına göre alfabetik olarak düzenlenmiştir. Birincil tarihsel metinlerde, metinde başvurulan baskı veya çeviri yılı esas alınmıştır.

\begingroup
\small
\newcommand{\refentry}[1]{\par\noindent\hangindent=1.5em\hangafter=1 #1\par\smallskip}

\refentry{Abrami, P. C., Bernard, R. M., Borokhovski, E., Waddington, D. I., Wade, C. A. ve Persson, T. (2015). Strategies for teaching students to think critically: A meta-analysis. \emph{Review of Educational Research, 85}(2), 275--314.}
\refentry{al-Khwārizmī, Muḥammad ibn Mūsā. (1831). \emph{The Algebra of Mohammed ben Musa} (F. Rosen, Çev.). Oriental Translation Fund.}
\refentry{Anellis, I. H. (2012). The genesis of the truth-table device. \emph{Logica Universalis, 6}, 87--110.}
\refentry{Appel, K. ve Haken, W. (1977). Every planar map is four colorable. Part I: Discharging. \emph{Illinois Journal of Mathematics, 21}, 429--490.}
\refentry{Appel, K., Haken, W. ve Koch, J. (1977). Every planar map is four colorable. Part II: Reducibility. \emph{Illinois Journal of Mathematics, 21}, 491--567.}
\refentry{Barroso, C., Ganley, C. M., McGraw, A. L., Geer, E. A., Hart, S. A. ve Daucourt, M. C. (2021). A meta-analysis of the relation between math anxiety and math achievement. \emph{Psychological Bulletin, 147}(2), 134--168.}
\refentry{Balacheff, N. (1988). Aspects of proof in pupils' practice of school mathematics. D. Pimm (Ed.), \emph{Mathematics, Teachers and Children} içinde (ss. 216--235). Hodder \& Stoughton.}
\refentry{Bashmakova, I. G. ve Smirnova, G. S. (2000). \emph{The Beginnings and Evolution of Algebra}. Mathematical Association of America.}
\refentry{Benacerraf, P. (1965). What numbers could not be. \emph{The Philosophical Review, 74}(1), 47--73.}
\refentry{Benacerraf, P. (1973). Mathematical truth. \emph{The Journal of Philosophy, 70}(19), 661--679.}
\refentry{Benacerraf, P. ve Putnam, H. (Ed.). (1983). \emph{Philosophy of Mathematics: Selected Readings} (2. bs.). Cambridge University Press.}
\refentry{Biggs, J. B. ve Collis, K. F. (1982). \emph{Evaluating the Quality of Learning: The SOLO Taxonomy}. Academic Press.}
\refentry{Boadu, S. K. ve Bonyah, E. (2024). The role of philosophy of mathematics education in mathematics teacher education. \emph{Cogent Education, 11}(1), 2433832. \url{https://doi.org/10.1080/2331186X.2024.2433832}}
\refentry{Boole, G. (1854). \emph{An Investigation of the Laws of Thought}. Walton and Maberly.}
\refentry{Bridges, D. ve Richman, F. (1987). \emph{Varieties of Constructive Mathematics}. Cambridge University Press.}
\refentry{Brouwer, L. E. J. (1913). Intuitionism and formalism. \emph{Bulletin of the American Mathematical Society, 20}, 81--96.}
\refentry{Bueno, O. (2020). Nominalism in the philosophy of mathematics. \emph{The Stanford Encyclopedia of Philosophy}.}
\refentry{Cajori, F. (1928). \emph{A History of Mathematical Notations}. Open Court.}
\refentry{Cantor, G. (1891). Ueber eine elementare Frage der Mannigfaltigkeitslehre. \emph{Jahresbericht der Deutschen Mathematiker-Vereinigung, 1}, 75--78.}
\refentry{Chace, A. B. (1927--1929). \emph{The Rhind Mathematical Papyrus} (2 cilt). Mathematical Association of America.}
\refentry{Chazan, D. (1993). High school geometry students' justification for their views of empirical evidence and mathematical proof. \emph{Educational Studies in Mathematics, 24}(4), 359--387.}
\refentry{Church, A. (1936). An unsolvable problem of elementary number theory. \emph{American Journal of Mathematics, 58}(2), 345--363.}
\refentry{Cohen, P. J. (1966). \emph{Set Theory and the Continuum Hypothesis}. W. A. Benjamin.}
\refentry{Colyvan, M. (2001). \emph{The Indispensability of Mathematics}. Oxford University Press.}
\refentry{Colyvan, M. (2012). \emph{An Introduction to the Philosophy of Mathematics}. Cambridge University Press.}
\refentry{De Millo, R. A., Lipton, R. J. ve Perlis, A. J. (1979). Social processes and proofs of theorems and programs. \emph{Communications of the ACM, 22}(5), 271--280.}
\refentry{De Toffoli, S. ve Mancosu, P. (2026). The philosophy of mathematical practice. \emph{The Stanford Encyclopedia of Philosophy}. \url{https://plato.stanford.edu/entries/mathematical-practice/}}
\refentry{De Villiers, M. (1990). The role and function of proof in mathematics. \emph{Pythagoras, 24}, 17--24.}
\refentry{Descartes, R. (1954). \emph{The Geometry of René Descartes} (D. E. Smith ve M. L. Latham, Çev.). Dover.}
\refentry{Detlefsen, M. (1986). \emph{Hilbert's Program: An Essay on Mathematical Instrumentalism}. Reidel.}
\refentry{Devlin, K. (2012). \emph{Introduction to Mathematical Thinking}. Keith Devlin.}
\refentry{Einstein, A. (1916). Die Grundlage der allgemeinen Relativitätstheorie. \emph{Annalen der Physik, 49}, 769--822.}
\refentry{Enderton, H. B. (2001). \emph{A Mathematical Introduction to Logic} (2. bs.). Academic Press.}
\refentry{Ernest, P. (1989). The impact of beliefs on the teaching of mathematics. P. Ernest (Ed.), \emph{Mathematics Teaching: The State of the Art} (ss. 249--254) içinde. Falmer Press.}
\refentry{Ernest, P. (1991). \emph{The Philosophy of Mathematics Education}. Falmer Press.}
\refentry{Euclid. (1956). \emph{The Thirteen Books of Euclid's Elements} (T. L. Heath, Çev.; 2. bs., 3 cilt). Dover.}
\refentry{Euler, L. (1988). \emph{Introduction to Analysis of the Infinite} (J. D. Blanton, Çev.). Springer.}
\refentry{Ewald, W. (2018). The emergence of first-order logic. \emph{The Stanford Encyclopedia of Philosophy}.}
\refentry{Field, H. (1980). \emph{Science Without Numbers: A Defence of Nominalism}. Princeton University Press.}
\refentry{Fischbein, E. (1982). Intuition and proof. \emph{For the Learning of Mathematics, 3}(2), 9--18, 24.}
\refentry{Fowler, D. (1999). \emph{The Mathematics of Plato's Academy: A New Reconstruction} (2. bs.). Oxford University Press.}
\refentry{Franklin, J. (2023). `Let no-one ignorant of geometry…': Mathematical parallels for understanding the objectivity of ethics. \emph{The Journal of Value Inquiry, 57}(2), 365--384.}
\refentry{Franzén, T. (2005). \emph{Gödel's Theorem: An Incomplete Guide to Its Use and Abuse}. A K Peters.}
\refentry{Frege, G. (1879). \emph{Begriffsschrift}. Louis Nebert.}
\refentry{Frege, G. (1997). \emph{The Frege Reader} (M. Beaney, Ed.). Blackwell.}
\refentry{Gentzen, G. (1935). Untersuchungen über das logische Schließen I. \emph{Mathematische Zeitschrift, 39}, 176--210.}
\refentry{Gonthier, G. (2008). Formal proof---the Four-Color Theorem. \emph{Notices of the American Mathematical Society, 55}(11), 1382--1393.}
\refentry{Gray, J. (2007). \emph{Worlds Out of Nothing: A Course in the History of Geometry in the 19th Century}. Springer.}
\refentry{Greenberg, M. J. (2008). \emph{Euclidean and Non-Euclidean Geometries: Development and History} (4. bs.). W. H. Freeman.}
\refentry{Gödel, K. (1940). \emph{The Consistency of the Axiom of Choice and of the Generalized Continuum-Hypothesis with the Axioms of Set Theory}. Princeton University Press.}
\refentry{Hales, T. C., Adams, M., Bauer, G., ve diğ. (2017). A formal proof of the Kepler conjecture. \emph{Forum of Mathematics, Pi, 5}, e2.}
\refentry{Hammack, R. (2018). \emph{Book of Proof} (3. bs.). Richard Hammack.}
\refentry{Hanna, G. (2000). Proof, explanation and exploration: An overview. \emph{Educational Studies in Mathematics, 44}, 5--23.}
\refentry{Harel, G. ve Sowder, L. (1998). Students' proof schemes: Results from exploratory studies. \emph{Research in Collegiate Mathematics Education III} içinde (ss. 234--283). American Mathematical Society.}
\refentry{Harper, R. (2016). \emph{Practical Foundations for Programming Languages} (2. bs.). Cambridge University Press.}
\refentry{Healy, L. ve Hoyles, C. (2000). A study of proof conceptions in algebra. \emph{Journal for Research in Mathematics Education, 31}(4), 396--428.}
\refentry{Heath, T. L. (1910). \emph{Diophantus of Alexandria: A Study in the History of Greek Algebra} (2. bs.). Cambridge University Press.}
\refentry{Heawood, P. J. (1890). Map-colour theorem. \emph{Quarterly Journal of Pure and Applied Mathematics, 24}, 332--338.}
\refentry{Heyting, A. (1956). \emph{Intuitionism: An Introduction}. North-Holland.}
\refentry{Hilbert, D. (1926). Über das Unendliche. \emph{Mathematische Annalen, 95}, 161--190.}
\refentry{Huffman, C. A. (2019). Pythagoras. \emph{The Stanford Encyclopedia of Philosophy}.}
\refentry{Iannone, P., Marasingha, G. ve Thoma, A. (2025). Teaching mathematics with Lean: Interactive theorem provers in the classroom. \emph{Notices of the American Mathematical Society, 72}(7), 740--747. \url{https://doi.org/10.1090/noti3184}}
\refentry{Kant, I. (1998). \emph{Critique of Pure Reason} (P. Guyer ve A. W. Wood, Çev.). Cambridge University Press.}
\refentry{Kempe, A. B. (1879). On the geographical problem of the four colours. \emph{American Journal of Mathematics, 2}(3), 193--220.}
\refentry{Knuth, E. J. (2002). Secondary school mathematics teachers' conceptions of proof. \emph{Journal for Research in Mathematics Education, 33}(5), 379--405.}
\refentry{Lakatos, I. (1976). \emph{Proofs and Refutations: The Logic of Mathematical Discovery}. Cambridge University Press.}
\refentry{Lear, J. (1982). Aristotle's philosophy of mathematics. \emph{The Philosophical Review, 91}(2), 161--192.}
\refentry{Linnebo, Ø. (2023). Platonism in the philosophy of mathematics. \emph{The Stanford Encyclopedia of Philosophy}. \url{https://plato.stanford.edu/entries/platonism-mathematics/}}
\refentry{Martinez, S. J. (2026). A Wittgensteinian account of the beauty of proofs. \emph{Synthese, 208}, 63. \url{https://doi.org/10.1007/s11229-026-05728-1}}
\refentry{Munkres, J. R. (2000). \emph{Topology} (2. bs.). Prentice Hall.}
\refentry{Netz, R. (1999). \emph{The Shaping of Deduction in Greek Mathematics}. Cambridge University Press.}
\refentry{Newton, I. (1999). \emph{The Principia: Mathematical Principles of Natural Philosophy} (I. B. Cohen ve A. Whitman, Çev.). University of California Press.}
\refentry{Oliveira e Silva, T., Herzog, S. ve Pardi, S. (2014). Empirical verification of the even Goldbach conjecture and computation of prime gaps up to $4\times10^{18}$. \emph{Mathematics of Computation, 83}(288), 2033--2060. \url{https://doi.org/10.1090/S0025-5718-2013-02787-1}}
\refentry{Peano, G. (1889). \emph{Arithmetices principia, nova methodo exposita}. Fratres Bocca.}
\refentry{Peano, G. (1897). Studii di logica matematica. \emph{Atti della Reale Accademia delle Scienze di Torino, 32}, 565--583.}
\refentry{Peirce, C. S. (1885). On the algebra of logic: A contribution to the philosophy of notation. \emph{American Journal of Mathematics, 7}(2), 180--202.}
\refentry{Peters, S. ve Westerståhl, D. (2006). \emph{Quantifiers in Language and Logic}. Oxford University Press.}
\refentry{Pierce, B. C. (2002). \emph{Types and Programming Languages}. MIT Press.}
\refentry{Popper, K. R. (2002). \emph{The Logic of Scientific Discovery}. Routledge.}
\refentry{Post, E. L. (1921). Introduction to a general theory of elementary propositions. \emph{American Journal of Mathematics, 43}(3), 163--185. \url{https://doi.org/10.2307/2370324}}
\refentry{Putnam, H. (1971). Philosophy of logic. J. P. De Cecco (Ed.), \emph{Mathematics, Matter and Method} içinde. Cambridge University Press.}
\refentry{Pólya, G. (2004). \emph{How to Solve It} (2. bs.). Princeton University Press.}
\refentry{Quine, W. V. O. (1948). On what there is. \emph{The Review of Metaphysics, 2}(5), 21--38.}
\refentry{Rav, Y. (1999). Why do we prove theorems? \emph{Philosophia Mathematica, 7}(1), 5--41. \url{https://doi.org/10.1093/philmat/7.1.5}}
\refentry{Reichenbach, H. (1938). \emph{Experience and Prediction}. University of Chicago Press.}
\refentry{Rittle-Johnson, B., Siegler, R. S. ve Alibali, M. W. (2001). Developing conceptual understanding and procedural skill in mathematics: An iterative process. \emph{Journal of Educational Psychology, 93}(2), 346--362. \url{https://doi.org/10.1037/0022-0663.93.2.346}}
\refentry{Rivest, R. L., Shamir, A. ve Adleman, L. (1978). A method for obtaining digital signatures and public-key cryptosystems. \emph{Communications of the ACM, 21}(2), 120--126. \url{https://doi.org/10.1145/359340.359342}}
\refentry{Rocq Development Team. (2025). \emph{Coq is now Rocq}. Rocq Prover. \url{https://rocq-prover.org/}}
\refentry{Russell, B. (1919). \emph{Introduction to Mathematical Philosophy}. George Allen \& Unwin.}
\refentry{Shannon, C. E. (1938). A symbolic analysis of relay and switching circuits. \emph{Transactions of the American Institute of Electrical Engineers, 57}(12), 713--723. \url{https://doi.org/10.1109/T-AIEE.1938.5057767}}
\refentry{Shapiro, S. (1997). \emph{Philosophy of Mathematics: Structure and Ontology}. Oxford University Press.}
\refentry{Shapiro, S. (2000). \emph{Thinking about Mathematics}. Oxford University Press.}
\refentry{Sheffer, H. M. (1913). A set of five independent postulates for Boolean algebras. \emph{Transactions of the American Mathematical Society, 14}(4), 481--488.}
\refentry{Skemp, R. R. (1976). Relational understanding and instrumental understanding. \emph{Mathematics Teaching, 77}, 20--26.}
\refentry{Solomon, M. (2001). Social empiricism. \emph{The Journal of Philosophy, 98}(9), 468--492.}
\refentry{Spinoza, B. (2002). \emph{Complete Works} (M. L. Morgan, Ed.; S. Shirley, Çev.). Hackett.}
\refentry{Stillwell, J. (2002). \emph{Mathematics and Its History} (2. bs.). Springer.}
\refentry{Stylianides, G. J., Stylianides, A. J. ve Moutsios-Rentzos, A. (2024). Research on proof and proving in mathematics education: A systematic review. \emph{ZDM--Mathematics Education, 56}, 47--59. \url{https://doi.org/10.1007/s11858-023-01518-y}}
\refentry{Tait, W. W. (1981). Finitism. \emph{The Journal of Philosophy, 78}(9), 524--546.}
\refentry{Taylor, R. ve Wiles, A. (1995). Ring-theoretic properties of certain Hecke algebras. \emph{Annals of Mathematics, 141}(3), 553--572. \url{https://doi.org/10.2307/2118560}}
\refentry{Thoma, A. ve Iannone, P. (2022). Learning about proof with the theorem prover Lean: The abundant numbers task. \emph{International Journal of Research in Undergraduate Mathematics Education, 8}, 64--93. \url{https://doi.org/10.1007/s40753-021-00140-1}}
\refentry{Toulmin, S. (1958). \emph{The Uses of Argument}. Cambridge University Press.}
\refentry{Turing, A. M. (1937). On computable numbers, with an application to the Entscheidungsproblem. \emph{Proceedings of the London Mathematical Society, s2-42}(1), 230--265. \url{https://doi.org/10.1112/plms/s2-42.1.230}}
\refentry{Tymoczko, T. (1979). The four-color problem and its philosophical significance. \emph{The Journal of Philosophy, 76}(2), 57--83. \url{https://doi.org/10.2307/2025976}}
\refentry{Usiskin, Z. (1988). Conceptions of school algebra and uses of variables. A. F. Coxford ve A. P. Shulte (Ed.), \emph{The Ideas of Algebra, K--12} içinde (ss. 8--19). National Council of Teachers of Mathematics.}
\refentry{Velleman, D. J. (2006). \emph{How to Prove It} (2. bs.). Cambridge University Press.}
\refentry{Viète, F. (1591). \emph{In artem analyticem isagoge}. Jamet Mettayer.}
\refentry{Wason, P. C. (1966). Reasoning. B. M. Foss (Ed.), \emph{New Horizons in Psychology} içinde (ss. 135--151). Penguin.}
\refentry{Weisstein, E. W. (2026). Goldbach conjecture. \emph{MathWorld---A Wolfram Web Resource}. \url{https://mathworld.wolfram.com/GoldbachConjecture.html}}
\refentry{Whitehead, A. N. ve Russell, B. (1910--1913). \emph{Principia Mathematica} (3 cilt). Cambridge University Press.}
\refentry{Wigner, E. P. (1960). The unreasonable effectiveness of mathematics in the natural sciences. \emph{Communications on Pure and Applied Mathematics, 13}(1), 1--14. \url{https://doi.org/10.1002/cpa.3160130102}}
\refentry{Wiles, A. (1995). Modular elliptic curves and Fermat's last theorem. \emph{Annals of Mathematics, 141}(3), 443--551. \url{https://doi.org/10.2307/2118559}}
\refentry{Yang, K., Poesia, G., He, J., Li, W., Lauter, K. E., Chaudhuri, S. ve Song, D. (2026). Formal reasoning meets LLMs: Toward AI for mathematics and verification. \emph{Communications of the ACM, 69}(3), 66--73. \url{https://doi.org/10.1145/3750038}}
\refentry{Zakariya, Y. F. (2022). Improving students' mathematics self-efficacy: A systematic review of intervention studies. \emph{Frontiers in Psychology, 13}, 986622. \url{https://doi.org/10.3389/fpsyg.2022.986622}}
\refentry{Zermelo, E. (1908). Untersuchungen über die Grundlagen der Mengenlehre I. \emph{Mathematische Annalen, 65}, 261--281. \url{https://doi.org/10.1007/BF01449999}}


\endgroup